% Pista alg-24j-q2 · Àlgebra
% Procedència: PAU Catalunya 2024 · Convocatòria de juny · Sèrie 1 · Qüestió 2
\documentclass[11pt,a4paper]{article}
\usepackage{pista}
\pistainfo{Catalunya 2024}{Convocatòria de juny}{Sèrie 1}

\begin{document}

\section*{\textcolor{blauPista}{Qüestió 2}}

\noindent Considereu el sistema d'equacions següent:
\[
\left.\begin{array}{r}
4x + 2y - z = 4 \\
x - y + kz = 3 \\
3x + 3y = 1
\end{array}\right\},
\]
on $k$ és un paràmetre real.

\subsection*{\textit{a)} \normalfont Discutiu el sistema per als diferents valors del paràmetre $k$, i resoleu-lo per a $k = 0$. \hspace*{\fill}\textnormal{[1~punt]}}

\pista{Aquest apartat té dues parts.

\medskip
\textit{Discutir.} Fes-ho amb el teorema de Rouché-Frobenius:
\begin{itemize}
\item escriu la matriu de coeficients $A$ i la matriu ampliada $A^{'}$;
\item calcula $\det(A)$ i troba els valors de $k$ que l'anul·len;
\item per a cadascun d'aquests valors, compara els rangs d'$A$ i d'$A^{'}$;
\item no oblidis el cas $\det(A) \neq 0$.
\end{itemize}
\textit{Resoldre per a $k = 0$.} Comprova primer si aquest valor anul·la $\det(A)$: què et diu això sobre el tipus de sistema i el nombre de solucions?}

\subsection*{\textit{b)} \normalfont Resoleu el sistema per a $k = -1$. \hspace*{\fill}\textnormal{[0,75~punts]}}

\pista{Recorda què has trobat a l'apartat \textit{a)} per a aquest valor de $k$.

\medskip
Si és compatible indeterminat, pren una variable com a paràmetre (per exemple $x = \lambda$) i expressa les altres dues en funció de $\lambda$.}

\subsection*{\textit{c)} \normalfont Per a $k = -1$, modifiqueu la tercera equació de manera que el sistema esdevingui incompatible. Justifiqueu la resposta. \hspace*{\fill}\textnormal{[0,75~punts]}}

\pista{Has d'aconseguir que els rangs d'$A$ i $A^{'}$ no coincideixin.
\begin{itemize}
\item Per a $k = -1$, fixa't quina relació hi ha entre la tercera equació i les altres dues.
\item Quina part de la tercera equació has de canviar perquè el rang d'$A$ es mantingui però el d'$A^{'}$ augmenti?
\item Justifica-ho comparant els rangs.
\end{itemize}}

\end{document}
